\documentclass[11pt,a4paper]{article}
\usepackage[utf8]{inputenc}
\usepackage[T1]{fontenc}
\usepackage[spanish]{babel}
\usepackage{amsmath,amssymb,mathtools}
\usepackage{geometry}
\usepackage{enumitem}
\usepackage{lmodern}
\usepackage{microtype}
\geometry{margin=2.4cm}
\setlength{\parindent}{0pt}
\setlength{\parskip}{0.45em}
\newcommand{\R}{\mathbb{R}}
\newcommand{\Ker}{\operatorname{Ker}}
\newcommand{\Ima}{\operatorname{Im}}
\newcommand{\rg}{\operatorname{rg}}
\newcommand{\tr}{\operatorname{tr}}
\newcommand{\Span}{\operatorname{L}}
\newcommand{\proy}{\operatorname{proy}}
\newcommand{\diag}{\operatorname{diag}}
\newcommand{\sg}{\operatorname{sg}}
\begin{document}

\begin{center}
{\LARGE\bfseries Álgebra Lineal}\\[0.4em]
{\Large Tema 6: El Espacio Vectorial Euclídeo --- Ejercicios}
\end{center}

\section{Ejercicio 1}
En $\R^3$ se considera:
\[
\langle(x,y,z),(x',y',z')\rangle
=
xx'+2yy'+2zz'+xz'+zx'+yz'+zy'.
\]
\begin{enumerate}[label=\alph*)]
\item Calcular su matriz en la base canónica.
\item Calcular la norma de $(1,1,1)$ y el ángulo entre $(1,0,1)$ y $(0,1,0)$.
\item Calcular:
\[
S^\perp,
\qquad
S=\Span\{(1,0,0),(0,1,0)\}.
\]
\end{enumerate}

\section{Ejercicio 2}
En $\R_2[x]$ sean:
\[
S=\{p(x):p(1)=0\},
\qquad
T=\{p(x):p(x)\text{ no tiene término independiente}\}.
\]
\begin{enumerate}[label=\alph*)]
\item Calcular una base de $S\cap T$ y otra de $S+T$.
\item Hallar, en la base canónica, unas ecuaciones implícitas de un subespacio suplementario de $S\cap T$.
\end{enumerate}

\section{Ejercicio 3}
En $\R^4$ con el producto escalar usual, obtener el subespacio ortogonal suplementario de:
\[
W=
\{(x,y,z,t):
x+y-z+t=0,\;
2x+y-z+3t=0\}.
\]

\section{Ejercicio 4}
Demostrar:
\[
\|u+v\|^2
=
\|u\|^2+\|v\|^2
\iff
\langle u,v\rangle=0.
\]

\section{Ejercicio 5}
Demostrar la ley del paralelogramo:
\[
\|u+v\|^2+\|u-v\|^2
=
2\left(\|u\|^2+\|v\|^2\right).
\]

\section{Ejercicio 6}
Demostrar:
\[
\|u\|=\|v\|
\iff
u+v\perp u-v.
\]
Deducir que un paralelogramo es un rombo si y sólo si sus diagonales son perpendiculares.

\section{Ejercicio 7}
En $\R^3$:
\[
\langle x,y\rangle
=
x_1y_1+2x_2y_2+3x_3y_3.
\]
Determinar una base de $F^\perp$, siendo:
\[
F=\{(x_1,x_2,x_3):x_1+x_2+2x_3=0\}.
\]

\section{Ejercicio 8}
En $\R^4$ con el producto escalar usual:
\begin{enumerate}[label=\alph*)]
\item Determinar un vector unitario ortogonal a:
\[
(1,2,1,0),\quad(0,-1,1,0),\quad(1,1,-2,1).
\]
\item Obtener por Gram--Schmidt una base ortonormal de:
\[
V=
\Span\{(1,2,-1,0),(0,1,1,0),(1,0,-2,1)\}.
\]
\end{enumerate}

\section{Ejercicio 9}
Calcular la proyección ortogonal de:
\[
(1,2,1)
\]
sobre:
\[
S=
\Span\{(0,1,2),(1,2,3)\}.
\]

\section{Ejercicio 10}
En $\mathcal P_1(x)$:
\[
\langle p,q\rangle
=
\int_0^1p(x)q(x)\,dx.
\]
\begin{enumerate}[label=\alph*)]
\item Obtener la matriz del producto escalar en la base canónica.
\item Calcular el ángulo entre $x+3$ y $2x+4$.
\item Calcular la proyección ortogonal de $x+3$ sobre $x+2$.
\item Determinar una base ortonormal a partir de la base canónica.
\end{enumerate}

\section{Ejercicio 11}
Obtener la proyección de:
\[
(1,1,1)
\]
sobre:
\[
S=
\Span\{(1,0,0),(0,1,0),(2,1,0)\}.
\]

\section{Ejercicio 12}
Dados:
\[
(1,1),\quad(2,4),\quad(3,7),\quad(4,9),
\]
ajustar por mínimos cuadrados:
\begin{enumerate}[label=\alph*)]
\item mediante un polinomio de grado $2$;
\item mediante un polinomio de grado $3$.
\end{enumerate}

\section{Ejercicio 13}
Se relacionan pesos $a_i$ de vehículos y precios $b_i$:
\[
\begin{array}{c|cccc}
a_i&0.8&1&1.2&1.3\\
\hline
b_i&1&2&3&5
\end{array}
\]
Hallar el mejor ajuste lineal por mínimos cuadrados.

\section{Ejercicio 14}
Para un cultivo se obtienen:
\[
\begin{array}{c|cccc}
t\text{ (h)}&1&2&3&4\\
\hline
N\text{ (miles)}&4&8&11&14
\end{array}
\]
Determinar la población aproximada al cabo de $7$ horas mediante un ajuste lineal.

\section{Ejercicio 15}
Determinar una matriz ortogonal $P$ que diagonalice ortogonalmente:
\[
A=
\begin{pmatrix}
1&1&0\\
1&2&1\\
0&1&1
\end{pmatrix}.
\]

\section{Ejercicio 16}
Sea $E$ un espacio euclídeo de dimensión $3$. La matriz del producto escalar en:
\[
B=\{e_1,e_1+e_2,e_1+e_2+e_3\}
\]
es:
\[
G=
\begin{pmatrix}
1&0&0\\
0&2&0\\
0&0&1
\end{pmatrix}.
\]
Sea:
\[
B'=\{e_1,e_2,e_3\}.
\]
Razonar si son ciertas o falsas:
\begin{enumerate}[label=\arabic*)]
\item
\[
\langle2e_1+e_2,2e_1+e_2\rangle=5.
\]
\item
\[
S_1=\Span\{e_1\},
\qquad
S_2=\Span\{e_1+e_2,e_3\}
\]
son ortogonales.
\item La matriz del producto escalar en $B'$ es:
\[
\begin{pmatrix}
1&2&0\\
2&5&1\\
0&1&2
\end{pmatrix}.
\]
\item El determinante de una matriz ortogonal es $1$ o $-1$.
\item Los únicos valores propios de una matriz ortogonal son $1$ o $-1$.
\end{enumerate}

\section{Ejercicio 17}
Demostrar:
\begin{enumerate}[label=\alph*)]
\item El producto de dos matrices ortogonales del mismo orden es ortogonal.
\item La inversa de una matriz ortogonal es ortogonal.
\end{enumerate}

\section{Ejercicio 18}
Calcular $s,t$ para que:
\[
A=
\frac17
\begin{pmatrix}
2&t&s\\
3&s&2\\
s&-2&t
\end{pmatrix}
\]
sea ortogonal.

\section{Ejercicio 19}
Demostrar que una matriz es ortogonal si y sólo si sus vectores columna forman un sistema ortonormal para el producto escalar usual.

\section{Ejercicio 20}
Diagonalizar ortogonalmente:
\[
A=
\begin{pmatrix}
1&0&2\\
0&-1&0\\
2&0&1
\end{pmatrix},
\]
obteniendo $P$ tal que:
\[
J=P^tAP.
\]

\section{Ejercicio 21}
Determinar si:
\[
A=
\begin{pmatrix}
3&3&1\\
3&3&1\\
1&1&3
\end{pmatrix},
\qquad
B=
\begin{pmatrix}
3&3&4\\
3&3&4\\
4&4&4
\end{pmatrix}
\]
son congruentes.

\section{Ejercicio 22}
Dada:
\[
M=
\begin{pmatrix}
1&-2\\
-2&1
\end{pmatrix},
\]
hallar una base ortonormal formada por vectores propios y la matriz diagonal correspondiente.

\section{Ejercicio 23}
Diagonalizar ortogonalmente:
\[
A=
\begin{pmatrix}
0&1&1\\
1&0&1\\
1&1&0
\end{pmatrix},
\]
obteniendo $P$ tal que:
\[
J=P^tAP.
\]

\section{Ejercicio 24}
Sea:
\[
G=
\begin{pmatrix}
2&-1&0\\
-1&1&\frac12\\
0&\frac12&2
\end{pmatrix}
\]
la matriz de Gram de un producto escalar. Obtener la proyección ortogonal de:
\[
(1,1,1)
\]
sobre:
\[
x+2z=0.
\]

\section{Ejercicio 25}
Dada:
\[
A=
\begin{pmatrix}
2&1&1\\
1&2&1\\
1&1&2
\end{pmatrix},
\]
\begin{enumerate}[label=\alph*)]
\item Comprobar que es diagonalizable en $\R$.
\item Comprobar la ortogonalidad de sus autovectores.
\item Encontrar una base ortonormal de autovectores.
\item Comprobar que la matriz de paso es ortogonal.
\item Clasificar:
\[
q(x)=x^tAx.
\]
\end{enumerate}

\section{Ejercicio 26}
Sea:
\[
f:\R^2\to\R^2
\]
definido por:
\[
f(1,1)=(-1,1),
\qquad
f(1,2)=(-1,2).
\]
Demostrar, en la base:
\[
B=\{(1,1),(1,2)\},
\]
que si $A=M_B(f)$ y $G_B$ es la matriz del producto escalar usual en $B$, entonces:
\[
G_B=A^tG_BA.
\]

\section{Ejercicio 27}
Sea $A$ una matriz simétrica $3\times3$ con traza:
\[
2a+b.
\]
El subespacio propio correspondiente a $a$ contiene:
\[
S=\{(x,y,z):x+y+z=0\}.
\]
\begin{enumerate}[label=\alph*)]
\item Calcular una forma diagonal $J$ y una matriz ortogonal $P$ tal que:
\[
J=P^tAP
\]
según los valores de $a$ y $b$.
\item Determinar para qué valores de $a,b$ puede ser $A$ la matriz de un producto escalar.
\item Para los valores naturales más pequeños que cumplen lo anterior, calcular el ángulo entre:
\[
(1,-1,0)
\quad\text{y}\quad
(1,1,1).
\]
\end{enumerate}

\section{Ejercicio 28}
Para:
\[
u=(1,1,1),
\]
\begin{enumerate}[label=\alph*)]
\item Calcular la matriz de Householder $H(u)$.
\item Obtener los transformados de:
\[
(3,3,3),\qquad
(1,-1,0),\qquad
(1,2,3).
\]
\end{enumerate}

\section{Ejercicio 29}
Sea $E$ un espacio euclídeo con base:
\[
B=\{e_1,e_2,e_3\}.
\]
Se sabe que:
\begin{itemize}
\item $e_2$ y $e_3$ forman un ángulo $\pi/3$;
\item $\|e_1\|,\|e_2\|,\|e_3\|$ son los números naturales pares más pequeños compatibles con el producto escalar;
\item
\[
B'=\{e_1,e_2+e_3,e_2-e_3\}
\]
es una base ortogonal.
\end{itemize}
Calcular:
\begin{enumerate}[label=\alph*)]
\item La matriz del producto escalar en $B$.
\item Una base ortonormal de:
\[
S^\perp,
\qquad
S=\Span\{e_1+e_2+e_3,e_1+2e_2\}.
\]
\item La proyección ortogonal de $e_2-e_3$ sobre $S^\perp$.
\end{enumerate}

\section{Ejercicio 30}
Dada la curva:
\[
C=
\{(x,y)\in\R^2:
10x^2-12xy+5y^2=1\},
\]
encontrar una base en la que la matriz asociada sea diagonal.

\section{Ejercicio 31}
En un espacio euclídeo $E$ de dimensión $3$, la matriz de Gram en:
\[
B=\{e_1,e_1+e_2,e_1+e_2+e_3\}
\]
es:
\[
G=
\begin{pmatrix}
1&0&0\\
0&2&0\\
0&0&1
\end{pmatrix}.
\]
Sea:
\[
B'=\{e_1,e_2,e_3\}.
\]
\begin{enumerate}[label=\alph*)]
\item Calcular:
\[
\langle2e_1+e_2,2e_1+e_2\rangle.
\]
\item Obtener la matriz del producto escalar en $B'$.
\item Comprobar si:
\[
V_1=\Span\{e_1\},
\qquad
V_2=\Span\{e_1+e_2,e_3\}
\]
son ortogonales.
\end{enumerate}

\end{document}
